\documentclass{article}

% Uses the NeurIPS template when supplied; otherwise a standard article.
\PassOptionsToPackage{numbers,compress}{natbib}
\IfFileExists{neurips_2025.sty}{%
  \usepackage[preprint]{neurips_2025}
}{%
  \usepackage[letterpaper,margin=1in]{geometry}
  \usepackage{natbib}
}
\usepackage[utf8]{inputenc}
\usepackage[T1]{fontenc}
\usepackage{graphicx}
\usepackage{float}
\usepackage{amsmath}
\usepackage{booktabs}
\usepackage{longtable}
\usepackage{microtype}
\usepackage{xcolor}
\usepackage{url}
\usepackage{array}
\usepackage{hyperref}
\hypersetup{colorlinks=true,linkcolor=blue,citecolor=blue,urlcolor=blue,
  pdftitle={Detecting AI Edits that Misrepresent Food Condition},pdfauthor={Mahad Faruqi}}
\bibliographystyle{abbrvnat}
\setlength{\emergencystretch}{2em}
\title{Detecting AI Edits that Misrepresent Food Condition}
\author{Mahad Faruqi \\
  Purdue EcoAI Lab \\
  Advisor: Prof. Haoran You}
\date{DoorDash AI Research Fellowship Proposal}

\begin{document}
\maketitle

\section{Motivation and Research Question}
Can we detect localized generative edits that make food appear burned, spoiled, or spilled, while avoiding false flags on genuine damage photos and harmless edits? I propose a controlled benchmark and a selective detection system for reviewing food-condition evidence. The central question is whether useful detection survives unfamiliar image editors, phone-camera variation, and image compression.

DoorDash's 2024 internship work used vision-transformer embeddings to identify reused refund photos.~\citep{doordashPhotoReuse2024} This proposal addresses a different case: a new photo whose food condition has been altered, even when no matching original is available. I would study the feasibility and limits of detecting such edits, without assuming their prevalence in DoorDash traffic or treating a suspicious image as proof of fraud.

My hypothesis is that local manipulation evidence, combined with the location of apparent food damage, can improve detection over a single whole-image AI score. The system must distinguish genuine visible damage from artificial changes to that damage. Its output would be a calibrated review signal with an uncertain option, not an automatic refund decision.

\section{Implementation Approach and Research Contribution}
\textbf{Construct matched counterfactual examples.} Start with approved food photos, including ordinary meals and genuine damage. Create controlled condition-changing edits alongside benign edits matched for editor, mask size, and location. Human reviewers would verify whether each edit actually changes perceived condition. Include difficult negatives such as naturally charred food, sauces, condensation, poor lighting, and compression artifacts. Retain source provenance, masks, prompts, model revisions, and seeds for reproducibility.

\textbf{Prevent shortcuts in the benchmark.} Apply the same resizing, re-encoding, and metadata handling to edited and real images. Keep all derivatives of a source in the same split. Hold out an editor family, merchants or capture sessions, and selected edit types. Use real damaged-food photos as a central negative class so the detector cannot succeed simply by learning that food which looks bad is suspicious. Source originals are available for annotation only; inference receives the submitted photo.

\textbf{Test a small, interpretable detector.} Reproduce a general AI-image detector, a local manipulation detector, and a vision-language judging baseline. Then test a lightweight verifier combining patch-level manipulation scores with food-region and apparent-damage localization. Train only the small adaptation or verification layers initially. Provide a suspect-region overlay, calibrated score, and abstention state; evaluate whether the highlighted region is correct rather than treating a plausible explanation as evidence.

Local-edit detection is an established research problem: DSRI evaluates detection and localization, and DailyBench covers modern synthesis and object-level manipulation.~\citep{dsriEditDetection2026,jiang2026dailybench} The proposed contribution is a food-condition benchmark and decision-focused evaluation that separates harmful content changes from benign processing, measures false flags on real complaints, and tests transfer to unseen editors.

\section{Evaluation Plan}
With an integrity research sponsor, define acceptable false-positive and review rates before tuning. Report detection recall at those operating points, precision across plausible prevalence levels, and confidence intervals. Size the genuine-photo test set to support the claimed false-positive ceiling; a small clean test set cannot establish a rare-error guarantee. Measure errors separately on genuine damaged food, undamaged food, and benignly edited images.

Evaluate localization against reviewed edit masks and report performance by edit size, food type, lighting, camera source, and postprocessing. Ablate local evidence, semantic localization, matched benign edits, and abstention. Compare methods under the same labeling and compute budgets. Reserve a blind test using unfamiliar editors and, if available, independently adjudicated operational cases. Synthetic accuracy alone would not justify deployment.

Measure p50/p95 inference latency, GPU memory, cost per reviewed image, and the fraction of cases escalated to people. Compare a lightweight first-stage model with a more expensive verifier called only on uncertain cases. This adds a systems question close to my thesis: when does selective computation preserve useful detection while reducing serving cost? All escalation costs and false flags count in the comparison.

Success would mean higher recall than strong baselines at the same genuine-photo false-positive and review budgets, with evidence of transfer to an unseen editor. If that transfer fails, the valuable outcome is a measured boundary on image-only detection and a recommendation to rely on additional evidence. The project would not infer customer intent, establish food safety from pixels, or automatically reject complaints.

\section{Preparation and Why DoorDash}
My master's research in Purdue's EcoAI Lab, advised by Prof. Haoran You, focuses on diffusion execution under memory and quality constraints. I have built reproducible runners and profiling infrastructure around PyTorch and stable-diffusion.cpp, analyzed A100 execution bottlenecks, and enabled quantized FLUX.2 [klein] generation on an 8 GB Jetson Orin Nano. These skills transfer to controlled image generation, auditable experiments, and inference-cost measurement. Image forensics would be a deliberate extension of my work, supported by collaboration with vision and integrity researchers.

DoorDash is a strong setting because real food photographs and review decisions expose failure cases that generic authenticity benchmarks miss. I would seek governed access to adjudicated cases, representative camera and upload conditions, and a sponsor who can define the cost of unnecessary customer review. Its earlier photo-reuse project provides a relevant operational precedent.~\citep{doordashPhotoReuse2024} The fellowship's evaluation and multimodal-understanding focus makes a rigorous reliability study valuable even if universal detection proves unattainable.

\section{Milestones and Feasibility}
\textbf{Weeks 1-2:} define the threat model, data permissions, and review targets. \textbf{Weeks 3-5:} build a pilot benchmark and reproduce baselines. \textbf{Weeks 6-8:} test local verification, calibration, and selective computation. \textbf{Weeks 9-12:} freeze models, run blind cross-editor tests, and deliver a benchmark specification, reviewer-facing prototype, and technical report. If operational data access is delayed, use licensed or consented images and clearly separate laboratory findings from production claims.

% References are embedded so this file compiles independently in the Codex editor.
% For a BibTeX project, replace the thebibliography environment below with
% \bibliography{references}; the accompanying references.bib retains cited keys.
\begin{thebibliography}{9}
\small
\bibitem{doordashPhotoReuse2024}
DoorDash. \emph{Part 1: DoorDash 2024 summer intern projects}. 2024. Section on photo recycling detection. \href{https://careersatdoordash.com/blog/part-1-doordash-2024-summer-intern-projects/}{Engineering blog}.

\bibitem{dsriEditDetection2026}
J.~Hostetler. \emph{The Image Edit Detection and Localization Challenge}. Digital Safety Research Institute, 2026. \href{https://dsri.org/blog/the-image-edit-detection-and-localization-challenge/}{Challenge report}.

\bibitem{jiang2026dailybench}
X.~Jiang et al. \emph{DailyBench: A Unified Benchmark for AI-Generated and Manipulated Images from Modern Generative Models}. 2026. \href{https://arxiv.org/abs/2607.24016}{arXiv:2607.24016}.

\end{thebibliography}
\end{document}
